\section{Conclusion}
\label{sec:conclusion}

We audited a complete prediction-to-policy air-quality pipeline on $1{,}180{,}551$
station-hours from $204$ Indian urban monitors and found that its most consequential
claims do not survive scrutiny. Empirically, spatio-temporal autocorrelation inflates
apparent predictive skill from an honest $R^2\approx0.22$--$0.30$ to a leaky $0.41$,
and the learned model is beaten by trivial persistence for temporal gap-filling;
pollutants add only $\Delta R^2\approx0.12$ beyond meteorology. Structurally, linear
causal discovery recovers a single robust edge---localized NO--O\textsubscript{3}
titration---rather than any regional photochemical mechanism, a ceiling imposed by the
structural absence of VOCs and sharpened by a verified ppb/\ugm{} unit-provenance
defect in three species. Algorithmically, we proved (\cref{lem:pareto}) that the true
Pareto frontier of the fitted linear surrogate is a closed-form line segment, against
which $100\%$ of the heuristic NSGA-II policies are dominated
($\mathrm{GD}=0.084$, $\mathrm{HVE}=6.6\%$).

Taken together these results motivate the four-part identification gate of
\cref{sec:gate} and a deliberate reframing. Rather than an over-claimed causal
source-apportionment and control-optimization system, the artifact is most honestly
and usefully released as an \emph{open, leakage-safe benchmark} for multi-pollutant
air-quality machine learning: a real national-scale dataset, a reproducible
evaluation protocol, a baseline ladder that future methods must beat, one
chemically interpretable hypothesis, and an explicit account of what each stronger
claim would require. All code, data, seeds and analysis scripts are released to this
end.

\paragraph{Future work.}
The productive path forward couples, rather than replaces, mechanism and data. We see
three directions: (i) \emph{hybrid modelling}, in which distributed observational
surrogates are coupled to, and constrained by, regional three-dimensional chemical
transport models (e.g.\ WRF-Chem, CMAQ) so that interventions are simulated through
\eqref{eq:adr} rather than fit on a correlational surface; (ii) \emph{speciation},
through low-cost VOC and carbonaceous-aerosol sensor arrays or satellite
HCHO/NO\textsubscript{2} proxies, to pass Gate~1 and enable genuine ozone-regime and
receptor-model analysis; and (iii) \emph{identified causal inference}, using documented
interventions or natural experiments with meteorological deconfounding and
pre-registered estimands. Until an analysis passes the identification gate, the
responsible contribution of observational environmental AI is a well-evaluated
predictive baseline and an honest map of its own limits.

\section*{Acknowledgements}
We thank the Central Pollution Control Board and the OpenAQ and Open-Meteo projects for
maintaining the open data infrastructure that made this audit possible.

\section*{Author contributions}
S.~Krishiykeishav, Pratyush Amit Sharma, Shreerag Venugopalan Nair, Jayshil Vinodbhai
Prajapati and Ryan Vinod jointly designed the study, implemented the pipeline and
audit, and wrote the manuscript.
